
\subsubsection{Interpretation of Perfect Correlation with High Divergence}

HumanEval and MBPP produce perfect model ranking correlation (ρ = 1.000) despite high distributional divergence (KL = 18.4). This pattern has a straightforward interpretation: benchmarks differ in difficulty calibration but measure the same underlying competency construct.

\textbf{Difficulty Differences:} MBPP shows consistently higher pass@1 scores (8-13 percentage points above HumanEval for the same models), indicating easier problems on average. Different score distributions do not imply different competencies being measured.

\textbf{Ranking Agreement:} Despite different absolute scores, models maintain identical rank positions. This indicates benchmarks agree on which models are stronger, even when they disagree on how much stronger.

\subsubsection{Competing Explanations}

\textbf{Explanation 1: Unidimensional Competency Space}

All execution-based code generation benchmarks measure general coding ability with different difficulty scaling. Perfect correlation suggests a single underlying dimension (general code generation competency) rather than multiple independent competencies (algorithmic clarity vs. practical patterns).

\textbf{Explanation 2: Sample Size Limitation}

Six models may be insufficient to detect ranking divergence. Perfect correlation could weaken with a larger, more diverse model population. However, ρ=1.0 is a strong signal unlikely to be solely due to sampling noise, and statistical significance (p<0.0001) confirms the pattern is non-random.

\textbf{Explanation 3: Benchmark Design Similarity}

HumanEval and MBPP may be more structurally similar than their design philosophies suggest. Both are Python-centric, function-level, unit-test-based generation tasks. True multi-dimensionality might emerge with cross-task-type comparisons (generation vs. understanding vs. repair).

We consider Explanation 1 most plausible given the strength of the correlation signal and consistency across different benchmark philosophies.

\subsubsection{Implications}

\textbf{Benchmark Selection:} If execution-based benchmarks rank models identically, evaluating on multiple such benchmarks provides redundant ranking information. One representative benchmark may suffice for ranking purposes.

\textbf{Aggregate Scoring:} Perfect correlation statistically justifies simple averaging across execution benchmarks as a composite score. Averaging would be inappropriate if benchmarks measured independent dimensions, but is valid when they measure the same construct.

\textbf{Evaluation Diversity:} Multi-dimensional evaluation likely requires task-type variety (generation + understanding + repair + translation) rather than multiple instances of the same task type (multiple generation benchmarks).

\subsubsection{Limitations}

\textbf{L1: Small Sample Size (n=6-8 models)}

Only 6 models have complete data on both benchmarks, below the originally planned 20+. This limits power to detect weak ranking divergence. However, perfect correlation is a strong signal that would unlikely reverse entirely with larger samples. The question is whether ρ=1.0 would weaken to ρ=0.85 (still high) or drop substantially.

\textbf{L2: Limited Benchmark Diversity (2 benchmarks)}

Analysis restricted to HumanEval and MBPP. APPS was unavailable due to dataset API changes. Two benchmarks suffice for testing pairwise correlation but prevent testing three-way factor structure. Results may not generalize beyond execution-based generation tasks.

\textbf{L3: Metric-Specific Analysis (pass@1 only)}

Rankings computed from pass@1 alone. Alternative metrics (runtime efficiency, error rates) might reveal different ordinal structures. However, pass@1 is the primary and most widely reported evaluation metric.

\textbf{L4: Uncontrolled Difficulty}

Analysis did not control for task difficulty. Distributional divergence may reflect difficulty differences rather than competency differences. Difficulty-adjusted analysis (item response theory) would separate difficulty variance from competency variance.

\textbf{L5: Incomplete Analysis Chain}

Only 2 of 5 planned sub-analyses were completed. Factor analysis, external validation, and intervention sensitivity were not tested. Conclusions are limited to data infrastructure (validated) and ranking correlation (perfect ρ observed).

\subsubsection{Broader Context}

This finding parallels debates in psychometric theory about general intelligence (g-factor) versus multiple intelligences. Our results suggest execution-based code evaluation may exhibit g-factor structure (unidimensional competency) rather than multiple-intelligences structure (task-specific competencies).

If models develop general coding representations that transfer broadly, this would explain why specialized benchmarks (algorithmic vs. practical) produce identical rankings. The alternative—that models develop independent competencies for different problem types—is not supported by the data.
